\section{MAP 把先验与似然合在一个点上}
\label{sec:11-Mathematical-Statistics/02-Point-Estimation/04}

上线前的最后一次调参，\texttt{weight\_decay} 从 \(10^{-4}\) 试到 \(10^{-2}\)，验证集 AUC 涨了 \(0.006\)，就定在 \(10^{-2}\)。评审会上有人问：为什么是 \(10^{-2}\)，不是 \(3\times 10^{-3}\)。标准答案是``网格搜出来的'', 这个答案不算错，只是它把一个可以说清楚的东西留成了黑箱。

其实那个数字有一句对应的人话。这份数据 \(n=10^{4}\)，残差标准差大约 \(0.5\)，取 \(\lambda=10^{-2}\) 等价于宣称：每个回归系数在看数据之前，被认为服从一个标准差 \(0.05\) 的零均值正态分布——也就是``我有九成把握每个系数落在 \(\pm 0.1\) 以内''。这句话可以拿去和业务方核对，网格搜索的编号不行。

把这个翻译关系讲清楚，是这一节的主要任务。它的出发点很朴素：在似然上乘一个先验，然后还是取峰。

\subsection{先验乘进来，峰就挪了}

频率学派把 \(\theta\) 当固定但未知的常数。Bayesian 换一个起点：看数据之前对 \(\theta\) 的判断本身写成一个分布 \(\pi(\theta)\)，看完数据之后用 Bayes 公式更新，得到后验

\[
\pi(\theta\given x)\propto L(\theta;x)\,\pi(\theta),
\]

更新机制本身见\hyperref[sec:05-Probability/02-Conditional-Probability/02]{``全概率公式与 Bayes 更新''}。后验是一个完整的答案，它对每个 \(\theta\) 都给了一个可信度。可交付物往往只要一个数，最省事的取法是取后验密度最高的那一点：

\[
\hat\theta_{\mathrm{MAP}}\in\argmax_\theta\ \bigl[\log L(\theta;x)+\log\pi(\theta)\bigr].
\]

比较一下：MLE 只有第一项，MAP 多了第二项。多出来的这一项不依赖数据，它是一股恒定的、把参数往先验偏好处拉的力。

拿一个能算到底的例子。抛硬币 \(n=10\) 次，\(S=7\) 次正面，MLE 是 \(0.7\)。先验取 \(\theta\sim\operatorname{Beta}(2,2)\)，一个温和地偏向公平的判断。Beta 与 Bernoulli 共轭，后验仍是 Beta，参数是 \(\operatorname{Beta}(\alpha+S,\ \beta+n-S)=\operatorname{Beta}(9,5)\)，见\hyperref[sec:13-Machine-Learning/06-Probabilistic-Models/07]{``共轭先验让 Bayesian 更新保持同一形式''}。对后验对数密度求导令零，得

\[
\hat\theta_{\mathrm{MAP}}=\frac{S+\alpha-1}{n+\alpha+\beta-2}=\frac{8}{12}=0.667 .
\]

条件对账：这个内部公式要求 \(\alpha+S>1\) 且 \(\beta+n-S>1\)，否则后验密度的峰在边界上，导数根本不过零。\(\operatorname{Beta}(0.5,0.5)\) 这类两端发散的先验就会触发这种情形，此时必须直接比较 \(\theta\to 0\) 与 \(\theta\to 1\) 处的极限，而不是套公式。

那个 \(\alpha-1\) 和 \(\beta-1\) 里的减一值得说一句。它意味着 \(\operatorname{Beta}(\alpha,\beta)\) 先验在 MAP 的意义下相当于事先见过 \(\alpha-1\) 次正面、\(\beta-1\) 次反面的虚拟数据，把它们和真实数据拌在一起再取比例。\hyperref[sec:09-Convex-Optimization/07-Applications/01]{``光滑经验风险与 \(L_2\) 正则化''}里 ridge 的增广写法做的是同一件事：往设计矩阵里塞 \(d\) 条虚拟样本，每条都说``这个系数应该是零''。两处的``虚拟数据''不是比喻，是同一个代数结构。

顺带留意，后验均值是 \((S+\alpha)/(n+\alpha+\beta)=9/14=0.643\)，和 MAP 的 \(0.667\) 不是一个数。众数是后验密度最高的点，均值是平方损失下的最优摘要，中位数是绝对损失下的最优摘要。三者各自对应一个损失函数，谁也不比谁更正宗，见\hyperref[sec:11-Mathematical-Statistics/07-Statistical-Decision-and-Reliable-Prediction/02]{``Bayes rule 最小化后验期望损失''}。

\subsection{负对数先验就是罚项}

把 MAP 的目标取负号写成最小化：

\[
\hat\theta=\argmin_\theta\ \underbrace{\bigl[-\log L(\theta;x)\bigr]}_{\text{数据项}}+\underbrace{\bigl[-\log\pi(\theta)\bigr]}_{\text{罚项}} .
\]

上一节说过，第一项就是训练用的负对数似然。第二项呢——它不依赖数据，只依赖参数，随着参数偏离先验偏好而增大。这正是正则项的定义。两者不是类比，是同一个表达式的两种叫法。

\begin{center}
\begin{tikzpicture}[font=\small,>=stealth]
  % ---- 左：先验密度 ----
  \node[font=\footnotesize,black!70] at (0,2.62) {先验密度 \(\pi(w)\)};
  \draw[->,black!60] (-2.35,0) -- (2.50,0) node[font=\footnotesize,black!70,right] {\(w\)};
  \draw[black!35,line width=0.9pt]
    plot[domain=-2:2,samples=80] (\x,{1.80*exp(-1.2*\x*\x)});
  \draw[black!85,line width=1.0pt]
    plot[domain=-2:2,samples=120] (\x,{1.80*exp(-1.55*abs(\x))});
  \node[font=\footnotesize,black!50,anchor=west] at (0.92,0.66) {Gauss};
  \node[font=\footnotesize,black!85,anchor=west] at (0.34,1.72) {Laplace};
  % ---- 中间箭头 ----
  \draw[->,black!65] (2.90,1.05) -- (3.90,1.05);
  \node[font=\footnotesize,black!60,anchor=south] at (3.40,1.12) {取 \(-\log\)};
  % ---- 右：罚项 ----
  \begin{scope}[shift={(6.70,0)}]
    \node[font=\footnotesize,black!70] at (0,2.62) {罚项 \(-\log\pi(w)\)};
    \draw[->,black!60] (-2.35,0) -- (2.50,0) node[font=\footnotesize,black!70,right] {\(w\)};
    \draw[black!35,line width=0.9pt]
      plot[domain=-1.98:1.98,samples=70] (\x,{0.56*\x*\x});
    \draw[black!85,line width=1.0pt] (-1.98,1.88) -- (0,0) -- (1.98,1.88);
    \node[font=\footnotesize,black!50,anchor=west] at (1.55,2.42) {\(\lambda\norm{w}_2^2\)};
    \node[font=\footnotesize,black!85,anchor=east] at (-2.05,2.10) {\(\lambda\norm{w}_1\)};
    \draw[black!85] (0,0) circle (3.4pt);
    \node[font=\footnotesize,black!60,anchor=north] at (0.05,-0.20) {尖角};
  \end{scope}
\end{tikzpicture}

\emph{左边是两个零均值先验的密度，右边是它们各自的负对数；取一次对数，先验就变成了教科书上的两个罚项。}
\end{center}

图上那条浅色曲线是零均值高斯先验，取负对数后是抛物线，正是 \(L_2\) 罚项。深色那条是 Laplace 先验，它在零点不可导，负对数是一个 V 形，正是 \(L_1\) 罚项。右图原点上圈出来的尖角不是画风问题——\hyperref[sec:09-Convex-Optimization/07-Applications/02]{``\(L_1\) 正则为何产生稀疏''}整篇讲的就是这个尖角，它让最优解有正概率恰好落在 \(w_j=0\) 上，而光滑的抛物线在原点导数为零，压不出精确的零。稀疏性的几何起点，在先验这一侧就是 Laplace 密度在零点的那个峰。

对应关系可以算到系数。线性模型 \(y=Xw+\varepsilon\)、\(\varepsilon\sim N(0,\sigma^2I)\)，先验 \(w\sim N(0,\tau^2I)\)，MAP 目标是

\[
\begin{aligned}
&\frac{1}{2\sigma^2}\norm{y-Xw}_2^2+\frac{1}{2\tau^2}\norm{w}_2^2\\
\propto\ &\frac{1}{2n}\norm{y-Xw}_2^2+\frac{\lambda}{2}\norm{w}_2^2,
\qquad \lambda=\frac{\sigma^2}{n\tau^2}.
\end{aligned}
\]

骨架一步：整体乘以 \(\sigma^2/n\)，不改变最优点。条件对账：用到了噪声同方差且各分量独立（否则 \(\sigma^2 I\) 要换成协方差矩阵，罚项变成 Mahalanobis 型）、先验各向同性且零均值（先验均值非零时罚项是 \(\norm{w-w_0}^2\)）。把高斯先验换成 \(\pi(w)\propto\prod_j e^{-\abs{w_j}/b}\)，同样的整理给出 \(\lambda=2\sigma^2/(nb)\)，罚项变成 \(\norm{w}_1\)，也就是 Lasso。

\subsection{调 \texorpdfstring{\(\lambda\)}{lambda} 就是在调先验的强度}

有了 \(\lambda=\sigma^2/(n\tau^2)\)，开头那个问题可以正着答也可以反着答。反解 \(\tau^2=\sigma^2/(n\lambda)\)，代入 \(\sigma=0.5\)、\(n=10^4\)、\(\lambda=10^{-2}\)，得 \(\tau^2=2.5\times 10^{-3}\)，\(\tau=0.05\)。所以那个网格点在说：每个系数的先验标准差是 \(0.05\)，\(95\%\) 的先验质量落在 \(\pm 0.1\) 之间。这句话业务方能听懂，也能反驳——如果某个特征的合理系数量级明明在 \(1\) 附近，那这个先验就是错的，该做的是给不同特征分组设不同的 \(\lambda\)，而不是继续加密网格。

这个公式还带出两条实践含义。分母上的 \(n\) 说明，先验强度固定时，最优的 \(\lambda\) 应该随样本量反比下降；数据变多，先验的相对话语权本来就该变小，把 \(\lambda\) 从小数据集原样搬到大数据集上，等于悄悄加强了先验。分子上的 \(\sigma^2\) 说明，噪声越大罚得越狠是合理的——信号弱的时候更该听先验的。这两条都和上一节那张似然曲线的图对得上：\(n\) 大则似然峰窄，先验拉不动它；\(\sigma\) 大则似然峰宽，先验一推就动。

至此可以回答一个常见的困惑：加正则到底是技巧还是原理。它是一句关于先验信念的声明，只不过通常没有被说出口。写下 \(\lambda\norm{w}_2^2\) 的人，无论自不自知，都已经断言了系数的先验分布形状和尺度。承认这一点的好处是这句声明变得可检验、可讨论、可以按特征分组定制，而不是一个只能靠验证集撞出来的魔数。

不过这条对应也不能推太远。同一个目标函数有三套讲法：控制复杂度、来自先验、改善条件数。三者指向同一个算式，对不确定性的后续陈述却不同。把 ridge 解当成``系数确实服从高斯先验''的推断结果，就要连带承认后验区间的那套语言；当成一个正则化技巧，那么区间得另找依据。算式相同，能说的话不同。

\subsection{MAP 只带走了后验的一个坐标}

后验是一整条分布，MAP 从中取走一个点。这个动作丢掉的东西比看上去多。

最直接的损失是不确定性。完整 Bayesian 对新样本的预测是后验预测分布 \(\int p(y\given\theta)\,\pi(\theta\given x)\,d\theta\)，它把参数的不确定性一起积分进去；MAP 用 \(p(y\given\hat\theta)\) 代替，等于假装参数已经被完全确定。小样本或高维时，这两者的预测区间宽度可以差出一大截，而 MAP 的那个更窄的区间会让系统显得比实际更自信。要保留这部分信息，路线是后验采样或变分近似，见\hyperref[sec:13-Machine-Learning/08-Sampling-and-Inference/03]{``变分推断把后验近似变成优化''}；区间形式的输出见\hyperref[sec:11-Mathematical-Statistics/04-Interval-Estimation/03]{``渐近区间与 Bayesian 可信区间''}。

更微妙的问题是众数在高维根本不代表典型情况。取 \(d\) 维标准正态后验，密度的最高点是原点，可几乎全部概率质量集中在半径约 \(\sqrt d\) 的一层薄壳上。\(d=1000\) 时，从后验里采一千个样本，没有一个会落在原点附近——众数所在的位置，是后验几乎从不去的地方。深度模型的参数量动辄百万，这个几何事实不是理论洁癖。MAP 依然有用，因为它可算，而完整后验往往不可算；但要清楚它交付的是``密度最高的点'', 不是``最典型的参数''。

还有一个纯技术的性质：MAP 不随参数化不变。密度在变量替换下带一个 Jacobian，

\[
\pi_\eta(\eta\given x)=\pi_\theta\bigl(g^{-1}(\eta)\given x\bigr)\left|\frac{d\theta}{d\eta}\right|,
\]

这个乘性因子一般会把峰挪走，所以 \(\hat\eta_{\mathrm{MAP}}\ne g(\hat\theta_{\mathrm{MAP}})\)。上一节提过 MLE 没有这个毛病，因为似然只比较同一个分布在不同标签下的取值，不涉及参数空间上的``单位体积''。用 \(\sigma\)、\(\log\sigma\) 还是 \(\sigma^2\) 做参数，MAP 会给出三个不同的答案。工程上的应对不复杂：把用了什么参数化写进文档，关键结论上顺手看一眼后验均值和后验区间是否指向同一个方向。

\subsection{先验从哪儿来，以及用了几次}

最后是一个容易滑过去的问题：\(\alpha,\beta\) 或 \(\tau^2\) 本身是谁定的。

凭领域知识定，没问题，那是货真价实的外部信息。但如果先用这批数据挑出一个``合适''的先验——比如让先验均值对上样本均值——再把它当成独立信息去算后验，同一批数据就被用了两遍：一遍定先验，一遍做更新。后验会比诚实的版本更窄，不确定性被系统性低估。经验 Bayes 有正规办法处理这种层次结构，把超参数自身的估计误差一并计入；随手看一眼数据再选个先验，不是经验 Bayes。

先验影响有多大取决于场景。固定维度、样本充足、条件正则时，似然项压倒温和先验，MAP 和 MLE 差不了多少。小样本、高维、参数接近不可识别时，先验可以主导结论。后一种情形下敏感性分析是必需动作：换一个同样合理的先验，方向是否翻转。翻转了，说明那个结论现在还是先验的结论，不是数据的结论。

这一单元的四种做法，判据、矩方程、似然、后验峰，都停在一个点上。可它们的表现好坏，从头到尾由同一样东西决定：数据里关于参数的信息有多少，以及似然的峰有多陡。那么这些信息能不能被压缩、压缩到什么程度还不损失，又该怎么度量？下一单元从这个问题开始。
